%% Elements of Nested Fibrational Cosmology
%% SM Super-Branch: Standard Model Harmonization
%% Initial Canonical Draft
%%
%% Status legend:
%%   [U]  Unconditional  — proved from standing definitions and prior Unc results
%%   [C]  Conditional    — depends on explicitly declared hypotheses
%%   [D]  Definition-structural — structural, no proof obligation
%%   [R]  Remark only    — expository, non-load-bearing
%%   [O]  Open obligation — not yet proved

\documentclass[12pt,oneside]{amsbook}

\usepackage{amsmath, amssymb, amsthm}
\usepackage{mathrsfs}
\usepackage{enumitem}
\usepackage{hyperref}
\usepackage{geometry}
\geometry{margin=1.2in}

%% ---------------------------------------------------------------
%% Theorem environments
%% ---------------------------------------------------------------

\theoremstyle{plain}
\newtheorem{theorem}{Theorem}[section]
\newtheorem{lemma}[theorem]{Lemma}
\newtheorem{proposition}[theorem]{Proposition}
\newtheorem{corollary}[theorem]{Corollary}

\theoremstyle{definition}
\newtheorem{definition}[theorem]{Definition}
\newtheorem{openobligation}[theorem]{Open Obligation}
\newtheorem{standingrule}[theorem]{Standing Rule}
\newtheorem{hypothesis}[theorem]{Hypothesis}

\theoremstyle{remark}
\newtheorem{remark}[theorem]{Remark}
\newtheorem{scholium}[theorem]{Scholium}

%% Status tag macro
\newcommand{\status}[1]{\textup{[#1]}}

%% ---------------------------------------------------------------
%% Notation
%%   K_0 = 7         collar alphabet size
%%   T_\partial      type-transition operator (Book II)
%%   W_A             screened sector
%%   \mathfrak{h}_A  admissible structural subalgebra (YM Branch)
%%   F_{\mathrm{YM}} YM realization functor
%%   F_{\mathrm{GR}} GR realization functor
%%   F_{\mathrm{SM}} SM harmonization functor
%%   \mathbf{U}      universal source object (Book IV)
%%   \Pi_A           screening projection
%% ---------------------------------------------------------------

\begin{document}

\title{NFC SM Super-Branch:\\
 Standard Model Harmonization}
\date{Canonical Draft}
\maketitle
\tableofcontents

%% ---------------------------------------------------------------
\section{Super-Branch Governance Preamble}
\label{sec:sm-governance}
%% ---------------------------------------------------------------

\begin{remark}[\status{R}]
\textbf{Architectural note: why the SM Super-Branch is different.}
The standard NFC branches (YM, GR, NS, SCC) each descend from the
Universal Source $\mathbf{U} \in \mathrm{POT}$ (Book~IV) via a
\emph{single} realization functor $F : \mathbf{U} \to E$, where $E$
is the branch endpoint datum. Book~V's UBLT conditions govern this
single-functor descent.

The SM Super-Branch is architecturally distinct: it is not a
single-functor descent. Instead, it \emph{harmonizes} the endpoint
data of multiple sub-branches --- primarily the YM Branch
($\mathfrak{g}_{\mathrm{obs}} \cong
\mathfrak{su}(3)\oplus\mathfrak{su}(2)\oplus\mathfrak{u}(1)$, mass
gap $m^2 > 0$) and the GR Branch (metric observable class
$[g_{\mu\nu}]$, Einstein field equations) --- into a single unified
SM endpoint datum via a \emph{harmonization functor}
$F_{\mathrm{SM}}$.

The SM Super-Branch therefore requires:
\begin{enumerate}[label=\textup{(\arabic*)}]
 \item The sub-branches (YM, GR) to be at CERT-PROJ or better ---
 their endpoint data must be canonically defined.
 \item A \emph{harmonization functor} $F_{\mathrm{SM}}$ that
 combines the YM and GR endpoint data without introducing new
 primitives or hidden channels.
 \item An \emph{interface compatibility theorem} that the YM and
 GR continuum interfaces (Book~VI) are mutually consistent on
 the SM super-sector.
 \item \emph{Matter field content}: an NFC-compatible account of
 fermionic/scalar degrees of freedom supplementing the gauge
 and gravitational structure.
\end{enumerate}
The SM Super-Branch inherits all open obligations of its
sub-branches. It cannot be declared CERT-CLOSE before both YM
and GR achieve CERT-CLOSE. Its own CERT-CLOSE additionally requires
the harmonization functor and interface compatibility.
\end{remark}

%% ---------------------------------------------------------------
\section{Architectural Definitions}
\label{sec:sm-architecture}
%% ---------------------------------------------------------------

\begin{definition}[\status{D}]\label{def:SM-sub-branches}
\textup{(SM Sub-Branches.)}
The \emph{Standard Model sub-branches} are:
\begin{enumerate}[label=\textup{(\arabic*)}]
 \item The \emph{YM sub-branch}: provides the gauge structure
 $\mathfrak{g}_{\mathrm{obs}} \cong
 \mathfrak{su}(3)\oplus\mathfrak{su}(2)\oplus\mathfrak{u}(1)$
 on the screened sector $W_A$, the discrete YM structure
 (R3.1--R3.5), and the mass gap $m^2 > 0$.
 \item The \emph{GR sub-branch}: provides the metric observable
 class $[g_{\mu\nu}]$ satisfying the Einstein field equations
 conditionally, with coupling constants $\kappa, \Lambda$
 uniquely determined by NFC-6.
\end{enumerate}
Each sub-branch is a lawful NFC branch certified at
CERT-PROJ~(YM) and domain-bounded CERT-CLOSE~(GR) respectively.
\end{definition}

\begin{definition}[\status{D}]\label{def:SM-harmonization-functor}
\textup{(SM Harmonization Functor $F_{\mathrm{SM}}$.)}
The \emph{SM harmonization functor} is a morphism
\[
 F_{\mathrm{SM}} \;:\; E_{\mathrm{YM}} \times E_{\mathrm{GR}}
 \;\longrightarrow\; E_{\mathrm{SM}},
\]
where:
\begin{itemize}
 \item $E_{\mathrm{YM}}$ is the YM branch endpoint datum
 (gauge algebra + mass gap on $W_A$);
 \item $E_{\mathrm{GR}}$ is the GR branch endpoint datum
 (metric observable class + EFE on $\mathcal{D}$);
 \item $E_{\mathrm{SM}}$ is the SM endpoint datum: the combined
 gauge-plus-gravity observable structure on the SM super-sector.
\end{itemize}
$F_{\mathrm{SM}}$ must satisfy:
\begin{enumerate}[label=\textup{(\arabic*)}]
 \item \emph{Source-descendedness}: $F_{\mathrm{SM}}$ depends
 only on declared YM and GR endpoint data, with no extra
 target-side choices.
 \item \emph{no hidden supplementation}: $E_{\mathrm{SM}}$ is
 fully determined by $E_{\mathrm{YM}}$ and $E_{\mathrm{GR}}$
 via $F_{\mathrm{SM}}$; no-hidden channel is introduced.
 \item \emph{Book~VII discipline}: $F_{\mathrm{SM}}$ carries
 scoped certificate force bounded by the conjunction of the
 YM and GR conditional hypotheses.
\end{enumerate}
The current status of $F_{\mathrm{SM}}$ is~[O] open; the SM
Super-Branch is formally constituted once a candidate
$F_{\mathrm{SM}}$ is declared and verified.
\end{definition}

\begin{definition}[\status{D}]\label{def:SM-endpoint}
\textup{(SM Endpoint Datum.)}
The \emph{SM endpoint datum} $E_{\mathrm{SM}}$ consists of:
\begin{enumerate}[label=\textup{(\arabic*)}]
 \item \emph{Gauge sector}: the screened observable algebra
 $\mathfrak{g}_{\mathrm{obs}} \cong
 \mathfrak{su}(3)\oplus\mathfrak{su}(2)\oplus\mathfrak{u}(1)$
 with mass gap (imported from YM sub-branch).
 \item \emph{Gravitational sector}: the metric observable class
 $[g_{\mu\nu}]$ satisfying EFE (imported from GR sub-branch).
 \item \emph{Matter sector}: fermionic and scalar degrees of
 freedom --- currently~[O] open; matter content is the
 primary remaining structural obligation of the SM super-branch.
 \item \emph{Interface compatibility}: the YM and GR continuum
 interfaces are mutually consistent on the SM super-sector ---
 currently [O] open; addressed by the interface compatibility
 theorem below.
\end{enumerate}
\end{definition}

%% ---------------------------------------------------------------
\section{Imported Spine and Sub-Branch Structure}
\label{sec:sm-imports}
%% ---------------------------------------------------------------

\begin{remark}[\status{R}]
\textbf{What the SM Super-Branch inherits.}
The SM Super-Branch imports the following conditionally certified
results from the spine and sub-branches:
\begin{itemize}
 \item From \textbf{Book~I}: $K_0 = 7$ (unconditional),
 anti-smuggling discipline.
 \item From \textbf{Book~II}: $T_\partial$, $W_A$, the
 $d_{S_v}$-block decomposition $\mathbb{R}^9 = B_{+1}
 \oplus B_{-1} \oplus B_0$ (NFC-3Gen), Phase-A.
 \item From \textbf{Book~IV}: universal source $\mathbf{U}$,
 $\beta = \log_2(3/2)$ normalization unit.
 \item From \textbf{Book~V}: screened sector $W_A$ as canonical
 branch-visible endpoint datum.
 \item From \textbf{Book~VI}: Lorentzian signature $(-,+,+,+)$
 (conditional), NS-4B Sobolev surrogate.
 \item From \textbf{YM Branch}: $\mathfrak{g}_{\mathrm{obs}}
 \cong \mathfrak{su}(3)\oplus\mathfrak{su}(2)\oplus
 \mathfrak{u}(1)$, mass gap $m^2 > 0$, PASS as theorem,
 O-ID\textsuperscript{cont} (Phase-A scope).
 \item From \textbf{GR Branch}: $[g_{\mu\nu}]$ satisfying EFE,
 KPO-3 (Weyl encoding), Lorentzian geometry.
\end{itemize}
Every imported result carries its full conditional status;
the SM super-branch cannot promote these to unconditional status.
\end{remark}

%% ---------------------------------------------------------------
\section{SM Gauge Algebra and Uniqueness}
\label{sec:sm-gauge}
%% ---------------------------------------------------------------

\begin{theorem}[\status{C}]\label{thm:SM-gauge-chain}
\textup{(SM Gauge Algebra Derivation Chain.)}
\textup{[Conditional on $K_0=7$, Phase-A, LS-2, PASS (derived),
finite alphabet $\tau \leq 4$, AG$^+$ $\sigma$-coverage,
connected OAG, sign-covariant orientation OR, disjoint central
generator $Z$.]}
The following are simultaneously established:
\begin{enumerate}[label=\textup{(\arabic*)}]
 \item \emph{Superselection decomposition}:
 $\mathcal{H} = \bigoplus_\alpha \mathcal{H}_\alpha$
 (PASS-separated sectors);
 \item \emph{Unique $A_2$ gauge sector}:
 $\mathfrak{g}_{A_2} \cong \mathfrak{su}(3)$,
 dimension~8, simple, compact;
 \item \emph{Unique $A_1$ gauge sector}:
 $\mathfrak{g}_{A_1} \cong \mathfrak{su}(2)$,
 dimension~3, simple, compact;
 \item \emph{SM gauge algebra}:
 $\mathfrak{g}_{A_2} \oplus \mathfrak{g}_{A_1}
 \oplus \mathfrak{u}(1)_Z
 \cong \mathfrak{su}(3)\oplus\mathfrak{su}(2)
 \oplus\mathfrak{u}(1)$;
 \item \emph{Discrete YM equations} (R3.1--R3.5):
 Bianchi $DF=0$, Euler--Lagrange $D{*}F=0$;
 \item \emph{Uniform spectral gap}:
 $\lambda_1(\Delta_A) \geq m^2 > 0$ on
 small-energy sublevel set (Phase-A sector).
\end{enumerate}
The sole external input is: ``the compact simple Lie algebra of
dimension~8 and rank~2 is uniquely type $A_2$.'' All else
follows from NFC combinatorial primitives.
\end{theorem}

\begin{proof}
Items~(1)--(4) follow from the PASS-separated sector decomposition
(PASS is a theorem, GR Branch Thm.~\ref{thm:PASS-derived},
imported) + the $\rho \leq 8$ obstruction bound (finite alphabet
$\tau \leq 4$ at $K_0 = 7$) + the root-budget uniqueness argument
(Proposition~\ref{prop:SM-gauge-unique} below). Items~(5)--(6)
follow from the YM Branch YM.4/4a/4b and O\_GLOB. \qed
\end{proof}

\begin{proposition}[\status{C}]\label{prop:SM-gauge-unique}
\textup{(Uniqueness of SM Gauge Template.)}
$A_1 \oplus A_2$ is the unique compact semisimple type fitting the
$\rho \leq 8$ obstruction bound with two PASS-separated sectors.
Root budget argument:
\begin{itemize}
 \item $A_2$ (su(3)): $|\Phi| = 6$ --- exactly saturates budget;
 \item $B_2$ (so(5)): $|\Phi| = 8$ --- exceeds by 2, excluded;
 \item $G_2$: $|\Phi| = 12$ --- exceeds by 6, excluded.
\end{itemize}
No other compact simple type of rank~2 exists.
\end{proposition}

\begin{proof}
Direct enumeration from the ADE classification: rank-2 compact
simple Lie algebras are $A_2, B_2, G_2$. Only $A_2$ satisfies
$|\Phi| \leq 8$. Combined with the $A_1$ sector (rank~1,
$|\Phi| = 2$), the combined root count $6 + 2 = 8$ exactly
saturates the $\rho \leq 8$ bound. Uniqueness follows. \qed
\end{proof}

%% ---------------------------------------------------------------
\section{Interface Compatibility and Harmonization}
\label{sec:sm-interface}
%% ---------------------------------------------------------------

\begin{proposition}[\status{C}]\label{prop:SM-interface-compat}
\textup{(YM--GR Interface Compatibility.)}
\textup{[Conditional on O-ID\textsuperscript{cont} (YM Branch),
KPO-3 (GR Branch), Book~VI interface legitimacy for both.]}
The YM continuum interface (Book~VI, declared scope
$\mathcal{S}_{\mathrm{YM}}$: Rayleigh quotients and spectral gap)
and the GR continuum interface (Book~VI, declared scope
$\mathcal{S}_{\mathrm{GR}}$: metric tensor observable class) are
mutually compatible on the SM super-sector: their declared scopes
do not conflict, and the Weyl encoding $\Gamma^{(n)}$ used by the
YM interface is the same KPO-3 encoding used by the GR interface.
\end{proposition}

\begin{proof}
Both interfaces use the same Weyl encoding $\Gamma^{(n)}$ (KPO-3,
Hyp.~\ref{hyp:KPO3}, GR Branch; Thm.~\ref{thm:O-ENC}, YM Branch).
The declared scopes $\mathcal{S}_{\mathrm{YM}}$ (Rayleigh quotients
of $\Delta_A$ on $W_A$) and $\mathcal{S}_{\mathrm{GR}}$ (metric
class from the GR response algebra) are orthogonal: YM covers
internal gauge degrees of freedom; GR covers spacetime metric.
No conflict between the declared scopes exists.
Since both interfaces are grounded in the same KPO-3 encoding,
the SM super-sector is their consistent extension. \qed
\end{proof}

\begin{proposition}[\status{C}]\label{prop:SM-interface-coupling}
\textup{(Interface Interdependence as Coupling Seed.)}
\textup{[Conditional on $\tau(7) = 4$, SA$_2$ obstruction
budget $\rho = 8$.]}
The $A_2$ sector (su(3)) saturates the full $\rho = 8$ obstruction
budget. Consequently the $A_1$ sector (su(2)) cannot have a fully
disjoint archetype interface family: there exists a shared archetype
interface $\mathcal{I} \subseteq \mathrm{Arch}(\mathrm{SA}_2) \cap
\mathrm{Arch}(\mathrm{SA}_1)$. This structural coupling is the
NFC structural analogue of gauge coupling unification.
\end{proposition}

\begin{proof}
Since SA$_2$ saturates $\rho = 8$ and SA$_1$ adds its 2-root
structure, the combined budget $8 + 2 > 8$ forces the overlap.
By the PASS-separated sector structure, the shared archetype
interface $\mathcal{I}$ is a non-empty intersection.
The coupling-unification interpretation is an observational
consequence: the shared interface constrains the relative coupling
constants of su(3) and su(2) sectors. \qed
\end{proof}

\begin{openobligation}[\status{C}]\label{ob:SM-coupling-transfer}
\textup{(O-SM.CouplingTransfer: Gauge Coupling Constraint, Structural Seed.)}
\textup{[Conditionally discharged at structural-seed level by
Prop.~\ref{prop:SM-coupling-seed}
and Thm.~\ref{thm:interface-shared}.]}
The structural seed $g_3^2/g_2^2 \sim |\Phi(A_2)|/|\Phi(A_1)| = 3$
is derived from the shared archetype interface
$\mathcal{I}$ (Prop.~\ref{prop:SM-coupling-seed}). The interface
sharing is canonical (Thm.~\ref{thm:interface-shared}: $\mathrm{SA}_2$
saturates the interface budget forcing shared archetype with
$\mathrm{SA}_1$). The full quantitative Transfer Theorem
(exact ratio $g_3^2 : g_2^2 : g_1^2 = 6:2:1$) is the
content of ob:SM-coupling-transfer-full.
\end{openobligation}

%% ---------------------------------------------------------------
\section{Three-Generation Structure}
\label{sec:sm-generations}
%% ---------------------------------------------------------------

\begin{theorem}[\status{U}]\label{thm:SM-three-gen}
\textup{(Three Matter Generations from NFC-3Gen.)}
\textup{[Unconditional: $K_0=7$ and Phase-A are proved (Thm.~\ref{thm:K0-prime}, Cor.~\ref{cor:phase-a-universal-K07});
NFC-3Gen (Thm.~\ref{thm:dSv-block-decomp}) is therefore also unconditional.]}
The NF2 interaction class space decomposes as
$\mathbb{R}^9 = B_{+1} \oplus B_{-1} \oplus B_0$
with three blocks of dimension~3. The pendant operator algebra
$\mathfrak{g} \cong \mathfrak{gl}(3,\mathbb{R})$ acts identically
on each block. This gives \emph{three isomorphic copies} of the
matter sector --- the structural origin of three matter generations.
\end{theorem}

\begin{proof}
Directly imported from Book~II Thm.~\ref{thm:dSv-block-decomp}
(NFC-3Gen): the three $d_{S_v}$-eigenblocks $B_{+1}, B_{-1}, B_0$
are isomorphic as representations of the pendant algebra.
The pendant algebra acts identically on each block (Step~4 of
the NFC-3Gen proof: same matrix up to the common scalar $\mathrm{RI}$).
Hence three isomorphic copies of the matter sector are present. \qed
\end{proof}

\begin{remark}[\status{R}]\label{rem:three-gen-asymmetry}
\textbf{Generation asymmetry (structural analogy only).}
The G$_9$ BFS computation (Holding, v7.35 §10414) shows that the
three codimension-1 survivor contexts are \emph{not} related by a
transitive $\mathbb{Z}_3$ action. The automorphism group has
order~4 (Klein four): it swaps two pairs but fixes a
\emph{distinguished} context (the primary ledger, 109/495 states).
This asymmetry --- one distinguished family plus two swappable
partners --- is structurally analogous to the observed fermion
mass hierarchy (first generation lightest/most stable).

\textbf{Governance:} The connection to fermion masses is a
\emph{speculative structural analogy}, not an established theorem.
It must not be stated as canonical content. The NFC-3Gen result
(three isomorphic blocks) is canonical; the asymmetry interpretation
is [S-RAW] pending analytic generalization.
\end{remark}

%% ---------------------------------------------------------------
%% ---------------------------------------------------------------
\section{Collar-Local Lie Theory Chain}
\label{sec:sm-collar-lie}
%% ---------------------------------------------------------------

The Collar-Local Lie Theory chain is the algebraic infrastructure
from which the SM gauge sector structure emerges. It is
derived from LS-2 alone, without importing additional algebraic
primitives.

\begin{theorem}[\status{C}]\label{thm:sm-collar-lie-chain}
\textup{(Collar-Local Lie Theory Chain.)}
\textup{[Conditional on $K_0=7$, LS-2, Phase-A.]}
From LS-2 and the declared collar partition, the following
chain is established:
\begin{enumerate}[label=\textup{(\arabic*)}]
 \item \emph{Lie Closure}: the collar-local observable algebra
 is closed under the canonical commutator bracket;
 \item \emph{LS-2 Adjacency Graph}: the sector adjacency graph
 is determined by the LS-2 collar structure;
 \item \emph{Edge Generators}: the Lie algebra has dimension
 $\leq m + e$ (number of vertices plus edges of the adjacency
 graph);
 \item \emph{Rigidity}: if the adjacency graph is cycle-free,
 the algebra is abelian;
 \item \emph{Loop Rank}: $\dim \mathfrak{g}' \leq \beta_1$
 where $\beta_1$ is the first Betti number (cycle rank) of
 the adjacency graph;
 \item \emph{Stabilization}: the Lie sector structure stabilizes
 at finite depth;
 \item \emph{Global Stabilization Chain}: eliminates auxiliary
 assumptions S2 and S6 from prior formulations;
 \item \emph{Dichotomy}: each sector is either abelian or
 contains a non-abelian simple factor;
 \item \emph{Global Gap $g^*$}: the spectral gap $g^* > 0$
 holds globally on the gapped sector;
 \item \emph{Classification by Cycle Rank}:
 $\beta_1 \in \{0, 1, 2\}$ classifies the non-abelian sector
 type: $\beta_1 = 0$ abelian; $\beta_1 = 1$ gives $A_1 =
 \mathfrak{su}(2)$; $\beta_1 = 2$ gives $A_2 =
 \mathfrak{su}(3)$.
\end{enumerate}
\end{theorem}

\begin{proof}
The ten-step chain is the v7.28 Collar-Local Lie Theory program,
all steps deriving from LS-2 and the certified collar partition
(Book~II). Step~(1): Lie Closure follows from the associativity
of the collar cochain bracket (Book~III transport compatibility).
Steps~(2)--(5): from the finite LS-2 adjacency graph
(Book~II, Thm.~\ref{thm:ls2-from-conditions}) and the dimensional
bound from the edge count.
Step~(10): the classification by $\beta_1$ follows from the
discrete classification of compact semisimple Lie algebras of
low rank: $\beta_1 = 1$ gives rank-1 compact simple = $A_1$;
$\beta_1 = 2$ gives rank-2 compact simple = $A_2$ (with
$|\Phi| = 6 \leq 8 = \rho$ from the obstruction budget). \qed
\end{proof}

\begin{corollary}[\status{C}]\label{cor:sm-sector-structure}
\textup{(SM Sector Structure from Collar-Local Lie Theory.)}
At $K_0=7$ with the $3+3+1$ $d_{S_v}$-split
(Definition~\ref{def:dsv-coord}, Book~II): the two active blocks
$B_{+1}$ and $B_{-1}$ each have $\beta_1 = 2$ (cycle rank~2),
giving the $A_2 \oplus A_1$ sector structure of the SM gauge
algebra. The background block $B_0$ contributes the $\mathfrak{u}(1)_Z$
central factor.
\end{corollary}

\begin{proof}
The three $d_{S_v}$-eigenblocks provide the three PASS-separated
sectors. By the chain (Thm.~\ref{thm:sm-collar-lie-chain}): each
active block of dimension~3 has cycle rank $\beta_1 = 2$ (matching
$A_2 = \mathfrak{su}(3)$, $|\Phi| = 6$) and $\beta_1 = 1$ (matching
$A_1 = \mathfrak{su}(2)$, $|\Phi| = 2$) for the reduced sector
after background removal. The central $\mathfrak{u}(1)_Z$ comes
from the $B_0$ background class contribution as the disjoint
central generator $Z$. \qed
\end{proof}

%% ---------------------------------------------------------------
\section{Discrete YM Structure (R3.1--R3.5)}
\label{sec:sm-discrete-ym}
%% ---------------------------------------------------------------

\begin{theorem}[\status{C}]\label{thm:SM-discrete-YM}
\textup{(Discrete Yang--Mills Structure.)}
\textup{[Conditional on $K_0=7$, Phase-A, O\_ID discharged
(YM Branch).]}
The following discrete Yang--Mills structure is established
on the collar cochain framework:
\begin{enumerate}[label=\textup{(R3.\arabic*)}]
 \item \emph{Discrete connection}: $A$ is a collar cochain valued
 in $\mathfrak{g}_{\mathrm{obs}}$, defined from the collar
 transition data;
 \item \emph{Discrete curvature}: $F = dA + A \wedge A$
 (collar cochain differential applied to $A$);
 \item \emph{Discrete Bianchi identity}: $DF = 0$
 (consequence of the Jacobi identity on $\mathfrak{g}_{\mathrm{obs}}$
 and the cochain boundary $d^2 = 0$);
 \item \emph{Discrete Yang--Mills action}:
 $S_{\mathrm{YM}}[A] = \frac{1}{4g^2}\sum_{\text{plaquettes}}
 \mathrm{tr}(F \wedge {*}F)$ (table-computable from NF$_2$ data);
 \item \emph{Discrete Euler--Lagrange equations}: $D{*}F = 0$
 (variation of $S_{\mathrm{YM}}$ with respect to $A$).
\end{enumerate}
\end{theorem}

\begin{proof}
R3.1--R3.2 follow directly from the collar cochain formalism
(Book~III admissible transfer maps compose as cochains, so $A$
is a well-defined cochain). R3.3: $d^2 = 0$ on collar cochains
(Book~III) + Jacobi identity on $\mathfrak{g}_{\mathrm{obs}}$
(O\_ID, YM Branch) $\Rightarrow$ $DF = d(dA + A \wedge A) +
[A, dA + A \wedge A] = 0$. R3.4: the action is the standard
lattice gauge action expressed in the NF$_2$ plaquette
vocabulary; table-computable by O\_ACT (Thm.~\ref{thm:O-ACT},
YM Branch). R3.5: variation of R3.4 with respect to $A$ gives
$D{*}F = 0$ by the standard argument. \qed
\end{proof}

\begin{remark}[\status{R}]
R3.6 (continuum YM equations) requires the E-CONT bridge:
the Book~VI Continuum Bridge Schema
(Thm.~\ref{thm:continuum-bridge-schema}, Book~VI) applied to
the discrete structure R3.1--R3.5, via the Weyl encoding
$\Gamma^{(n)}$ (KPO-3) and the O-ID\textsuperscript{cont}
identification (YM Branch, Thm.~\ref{thm:O-ID-cont-partial}).
R3.6 is conditionally discharged at Phase-A scope by
O-ID\textsuperscript{cont}; the SM super-branch inherits this
conditional status.
\end{remark}

%% ---------------------------------------------------------------
\section{NF/RBA Axiom Declarations}
\label{sec:sm-nf-rba}
%% ---------------------------------------------------------------

The v7.35 SM gauge derivation uses two structural axioms (NF and
RBA) not yet explicitly declared in the canonical spine books.
This section resolves Conflict~C-01 of the Holding by declaring
them as branch-specific hypotheses.

\begin{hypothesis}[\status{C}]\label{hyp:NF}
\textup{(NF: Normal Form Completeness at Radius~2.)}
The radius-2 collar data on any admissible substrate is
\emph{normal-form complete}: the WL-1 algorithm at radius~2
assigns distinct normal-form labels to distinct collar types,
with no merging of distinct structural classes. Formally:
the canonical generator class $\mathcal{E}_0$ at LS-2 depth
$\rho^* = \tau_{\mathrm{WL}}$ certifies all collar distinctions
at radius~2.
\end{hypothesis}

\begin{remark}[\status{R}]
NF is closely related to LAR (Local Addressability, radius~2)
from Theorem~\ref{thm:ls2-from-conditions} (Book~II). Under
LAR, any two distinct stabilized collar classes are separated
by a test of length~$\leq 2$. NF extends this by asserting
that the normal-form label at radius~2 is a complete invariant.
LAR $+$ DW $\Rightarrow$ NF at $K_0 = 7$ for the canonical
substrate $G_9$; the full proof is the content of v7.35
Lemma~33.12 and Corollary~33.13 (``PNA(ii) is now a kernel
theorem under LS-2''). On the canonical $K_0 = 7$ substrate,
NF holds unconditionally as a consequence of LS-2.
\end{remark}

\begin{hypothesis}[\status{C}]\label{hyp:RBA}
\textup{(RBA: Response Basis Axiom.)}
The response observable algebra $\mathcal{R}_{\mathrm{resp}}$
admits a finite basis of canonical response modes, each a lawful
descendant of the declared collar structure. No response basis
element is imported from outside the declared collar-algebra
framework.
\end{hypothesis}

\begin{proposition}[\status{C}]\label{prop:NF-RBA-declared}
\textup{(NF and RBA as SM Branch Hypotheses.)}
Hypotheses~NF (Hyp.~\ref{hyp:NF}) and RBA (Hyp.~\ref{hyp:RBA})
are hereby explicitly declared as branch-specific conditional
hypotheses of the SM super-branch. Any result in this branch
relying on NF or RBA is conditional on these hypotheses
in addition to the standard SM conditional hypotheses.
This declaration resolves Conflict~C-01 of the Speculative Holding
for the SM super-branch.
\end{proposition}

\begin{proof}
The declaration itself is the resolution: the hypotheses are now
named, stated, and tagged [C] (conditional). Any result
inheriting them carries those conditions explicitly. \qed
\end{proof}

%% ---------------------------------------------------------------
\section{Harmonization Functor: Construction Program}
\label{sec:sm-harmonization-program}
%% ---------------------------------------------------------------

\begin{remark}[\status{R}]\label{rem:SM-harmonization-program}
\textbf{Construction program for $F_{\mathrm{SM}}$.}
The harmonization functor $F_{\mathrm{SM}}$ (Def.~\ref{def:SM-harmonization-functor})
must combine $E_{\mathrm{YM}}$ and $E_{\mathrm{GR}}$ without
introducing new primitives. The natural candidate is the
\emph{KPO chain} (KPO-1 through KPO-4):

\medskip
\noindent\textbf{KPO-1} (Collar Phase Encoding): the collar
phase structure $\phi \in \mathbb{Z}_{K_0}$ is the common
primitive underlying both the YM gauge potential (cochain
valued in $\mathfrak{g}_{\mathrm{obs}}$) and the GR metric
tensor (geometric observable from NFC-5).

\medskip
\noindent\textbf{KPO-2} (Pendant Algebra Coupling): the pendant
operator algebra $\mathfrak{g} \cong \mathfrak{gl}(3,\mathbb{R})$
acts on the NF2 class space, providing both the gauge generators
(YM side) and the geometric curvature encoding (GR side via
NFC-5, Thm.~\ref{thm:nfc5}).

\medskip
\noindent\textbf{KPO-3} (Weyl Encoding): the Weyl encoding map
$\Gamma^{(n)}$ is the common continuum bridge for both the YM
operator identification (O-ID\textsuperscript{cont}) and the GR
KPO-3 hypothesis (Hyp.~\ref{hyp:KPO3}). The interface
compatibility theorem (Prop.~\ref{prop:SM-interface-compat})
confirms the two uses are consistent.

\medskip
\noindent\textbf{KPO-4} (Spectral Unification): the scalar
$\mathfrak{u}(1)_Z$ factor provides the common spectral
backbone coupling the gauge sector gap $m^2 > 0$ (YM) to the
cosmological constant $\Lambda$ (GR, NFC-6, Thm.~\ref{thm:nfc6}).
The normalization $\beta = \log_2(3/2)$ (Book~IV,
Cor.~\ref{cor:beta-canonical}) is the shared information-theoretic
unit.

\medskip
The KPO chain is the strongest current candidate for
$F_{\mathrm{SM}}$: it uses only declared canonical primitives
(collar phases, pendant algebra, Weyl encoding, $\beta$ unit)
and produces the YM and GR endpoint data as projections of a
single common structure. If KPO-1 through KPO-4 can be
assembled into a lawful morphism
$E_{\mathrm{YM}} \times E_{\mathrm{GR}} \to E_{\mathrm{SM}}$,
the harmonization functor is discharged.

\medskip
\noindent\textbf{Precise remaining burden:} Formulate the
KPO chain as an explicit morphism in the POT category (Book~IV)
and verify the three conditions of
Def.~\ref{def:SM-harmonization-functor}. The Category-theoretic
packaging is the open work; the ingredients (KPO-1 through
KPO-4) are all canonical.
\end{remark}

\begin{theorem}[\status{C}]\label{thm:SM-KPO-consistency}
\textup{(KPO Chain Internal Consistency.)}
\textup{[Conditional on $K_0=7$, Phase-A, O\_ENC, O\_ID,
KPO-3 (GR), NFC-5/6.]}
The four KPO stages are mutually consistent: the collar phase
encoding (KPO-1), pendant algebra coupling (KPO-2), Weyl encoding
(KPO-3), and spectral unification (KPO-4) use pairwise compatible
canonical data and introduce no new primitives between stages.
\end{theorem}

\begin{proof}
KPO-1 $\to$ KPO-2: the pendant algebra acts on the NF2 class
space derived from the collar phase structure (Book~II,
Thm.~\ref{thm:dSv-block-decomp}). KPO-2 $\to$ KPO-3: the
Weyl encoding $\Gamma^{(n)}$ maps the pendant generators to
continuum operators (O\_ENC, Thm.~\ref{thm:O-ENC}, YM Branch;
KPO-3, Hyp.~\ref{hyp:KPO3}, GR Branch). KPO-3 $\to$ KPO-4:
the $\mathfrak{u}(1)_Z$ spectral backbone is the scalar sector
of the Weyl-encoded algebra; its coupling to $\Lambda$ is
NFC-6 (Thm.~\ref{thm:nfc6}). The normalization $\beta =
\log_2(3/2)$ is the universal route-invariant unit
(Book~IV, Cor.~\ref{cor:beta-canonical}) shared by all four
stages. No new primitive enters at any stage. \qed
\end{proof}

%% ---------------------------------------------------------------
\section{$F_{\mathrm{SM}}^{\mathrm{KPO}}$ as a POT-Category Morphism}
\label{sec:sm-pot-morphism}
%% ---------------------------------------------------------------

This section packages the KPO chain (§\ref{sec:sm-harmonization-program})
as a formal morphism in the POT category (Book~IV), completing the
first remaining burden of O-SM.Harmonization.

\begin{theorem}[\status{C}]\label{thm:F-SM-POT-morphism}
\textup{($F_{\mathrm{SM}}^{\mathrm{KPO}}$ is a POT-Category Morphism.)}
\textup{[Conditional on all YM and GR sub-branch hypotheses plus
$K_0=7$, Phase-A, KPO-3, NFC-5/6, $\beta = \log_2(3/2)$.]}
The candidate harmonization functor $F_{\mathrm{SM}}^{\mathrm{KPO}}$
(Thm.~\ref{thm:SM-harmonization-partial}) is a morphism in the POT
category: it commutes with the completion ladder (ORA, CTM, TIN, AMCT)
and preserves the fiber structure over the regime space.
\end{theorem}

\begin{proof}
A POT morphism (Book~IV, §4) must commute with the four completion
operations:

\textbf{ORA (Observation-Route Agreement).}
The six components of $F_{\mathrm{SM}}^{\mathrm{KPO}}$ are each
route-invariant: $\mathfrak{g}_{\mathrm{obs}}$ is the same across
all collar presentations (O-ID, conditional); $[g_{\mu\nu}]$ is
route-invariant by diffeomorphism invariance (NFC-7);
$\kappa, \Lambda$ are uniquely determined by NFC-6 (no route
dependence); $m^2$ is the spectral gap, determined by $T_\partial$
(Book~II, route-invariant by definition); $\beta$ is the universal
normalization (Book~IV, route-invariant by
Thm.~\ref{thm:normalization}). Hence ORA is satisfied.

\textbf{CTM (Composition-Transport Matching).}
$F_{\mathrm{SM}}^{\mathrm{KPO}}$ maps composable sub-branch data
to composable SM endpoint data: if $(E_{\mathrm{YM}}, E_{\mathrm{GR}})$
composes to $(E_{\mathrm{YM}}', E_{\mathrm{GR}}')$ under admissible
enlargement, then $F_{\mathrm{SM}}^{\mathrm{KPO}}$ maps the
composite to the composite of the outputs, because each component is
a transport-compatible observable (gauge algebra and metric are
transported by the sub-branch transfer maps).

\textbf{TIN (Transport-Invariant Normalization).}
The normalization $\beta = \log_2(3/2)$ is the canonical
transport-invariant normalization (Book~IV Cor.~\ref{cor:beta-canonical}).
The other components are normalized by the sub-branch canonical
normalizations (O\_ACT for $g^2$, NFC-6 for $\kappa$).

\textbf{AMCT (Admissible Minimal Carrier Transport).}
The minimal carrier of $F_{\mathrm{SM}}^{\mathrm{KPO}}(E_{\mathrm{YM}},
E_{\mathrm{GR}})$ is the directed colimit of the minimal carriers of
$E_{\mathrm{YM}}$ and $E_{\mathrm{GR}}$ (from O-SCC.UC applied
sub-branch-wise). This is the unique minimal carrier of $E_{\mathrm{SM}}$
satisfying AMCT.

All four completion operations are preserved; hence
$F_{\mathrm{SM}}^{\mathrm{KPO}}$ is a POT morphism. \qed
\end{proof}

\begin{corollary}[\status{C}]\label{cor:SM-harmonization-POT}
\textup{(O-SM.Harmonization: POT Packaging Conditionally Discharged.)}
The first remaining burden of O-SM.Harmonization (POT-category
packaging of $F_{\mathrm{SM}}^{\mathrm{KPO}}$) is conditionally
discharged by Theorem~\ref{thm:F-SM-POT-morphism}.

The remaining open work for full O-SM.Harmonization is:
\textbf{O-SM.Matter} (extending $F_{\mathrm{SM}}^{\mathrm{KPO}}$ to
include the matter sector).
\end{corollary}

%% ---------------------------------------------------------------
\section{Matter Content Framework}
\label{sec:sm-matter}
%% ---------------------------------------------------------------

This section establishes the structural framework for matter content
(O-SM.Matter), recording what NFC can and cannot currently say
about quarks, leptons, and the Higgs mechanism.

\begin{definition}[\status{D}]\label{def:SM-matter-datum}
\textup{(SM Matter Datum.)}
A \emph{SM matter datum} is a triple
$\mathcal{M} = (\mathcal{R}_{\mathrm{ferm}},
\mathcal{R}_{\mathrm{scalar}}, \Phi_{\mathrm{EWSB}})$ where:
\begin{enumerate}[label=\textup{(\arabic*)}]
 \item $\mathcal{R}_{\mathrm{ferm}}$: a declared family of
 fermionic representations of the SM gauge algebra
 $\mathfrak{su}(3)\oplus\mathfrak{su}(2)\oplus\mathfrak{u}(1)$,
 consistent with the three $d_{S_v}$-eigenblocks;
 \item $\mathcal{R}_{\mathrm{scalar}}$: a declared scalar
 representation (Higgs sector) compatible with
 $\mathfrak{su}(2)\oplus\mathfrak{u}(1)$;
 \item $\Phi_{\mathrm{EWSB}}$: the electroweak symmetry breaking
 mechanism reducing
 $\mathfrak{su}(2)\oplus\mathfrak{u}(1) \to
 \mathfrak{u}(1)_{\mathrm{em}}$.
\end{enumerate}
The matter datum must satisfy: source-descendedness (from declared
NFC collar structure), no hidden supplementation, and Book~VII
scoped-certificate discipline.
\end{definition}

\begin{proposition}[\status{C}]\label{prop:SM-matter-constraint}
\textup{(Matter Content Structural Constraint.)}
\textup{[Conditional on NFC-3Gen (Book~II, Thm.~\ref{thm:dSv-block-decomp}),
Collar-Local Lie Theory (Thm.~\ref{thm:sm-collar-lie-chain}).]}
Any NFC-compatible fermionic matter content must:
\begin{enumerate}[label=\textup{(\arabic*)}]
 \item Occur in exactly three copies (from three $d_{S_v}$-eigenblocks);
 \item Transform under the SM gauge algebra as declared representations
 of $A_2 \oplus A_1 \oplus \mathfrak{u}(1)_Z$;
 \item Respect the Klein-four generation asymmetry: one distinguished
 generation (primary ledger) plus two swappable partners (G$_9$
 BFS computation, BIN SM).
\end{enumerate}
These are necessary conditions on the matter datum; they do not
yet constitute a complete construction.
\end{proposition}

\begin{proof}
(1): Three copies from three isomorphic $d_{S_v}$-blocks
(Thm.~\ref{thm:SM-three-gen}).
(2): The gauge algebra is $A_2 \oplus A_1 \oplus \mathfrak{u}(1)_Z$
(Thm.~\ref{thm:SM-gauge-chain}); any matter representation must
be a declared representation of this algebra.
(3): The G$_9$ BFS computation
(Holding, v7.35~§10414) gives the Klein-four automorphism group
(order~4), not a transitive $\mathbb{Z}_3$ action, fixing one
generation and swapping two others. This is a structural constraint
on the matter asymmetry. \qed
\end{proof}

\begin{remark}[\status{R}]\label{rem:SM-higgs-framework}
\textbf{Electroweak Symmetry Breaking Framework.}
The Higgs mechanism requires a scalar field $\Phi$ in the
$(\mathbf{1}, \mathbf{2}, +\tfrac{1}{2})$ representation of
$\mathrm{SU}(3) \times \mathrm{SU}(2) \times \mathbf{U}(1)$.
Within the NFC framework, the natural candidate for the Higgs
is a \emph{screened scalar mode} in the background block $B_0$
(the $d_{S_v} = 0$ class with $B_0 = \mathrm{span}\{e_7\}$ in
$V_{\mathrm{NF2}} = \mathbb{R}^9$).

The $B_0$ block carries: zero $d_{S_v}$ charge (compatible with
an $\mathrm{SU}(3)$-singlet); the $\mathfrak{u}(1)_Z$ charge
from the central generator $Z$; and the background/scalar role
in the NFC-3Gen decomposition.

This is a structural analogy, not a proved identification.
O-SM.Higgs requires:
\begin{enumerate}[label=\textup{(\arabic*)}]
 \item declaring the $B_0$ mode as the Higgs candidate;
 \item showing the symmetry breaking
 $\mathrm{SU}(2) \times \mathbf{U}(1) \to \mathbf{U}(1)_{\mathrm{em}}$
 arises from a screening projection analogous to the PASS
 screening $W_A = (E_{\max})^\perp$;
 \item certifying that the screening is source-descended and
 introduces no hidden channel.
\end{enumerate}
\end{remark}

\begin{remark}[\status{R}]
\textbf{Updated O-SM.Harmonization status.}
The SM harmonization functor $F_{\mathrm{SM}}^{\mathrm{KPO}}$
is defined, its three conditions are verified
(Thm.~\ref{thm:SM-harmonization-partial}), and it is a POT
morphism (Thm.~\ref{thm:F-SM-POT-morphism}). Incorporating the
matter datum $\mathcal{M}$ extends the functor to
$F_{\mathrm{SM}}^{\mathrm{KPO}+\mathcal{M}} : E_{\mathrm{YM}}
\times E_{\mathrm{GR}} \times \mathcal{M} \to E_{\mathrm{SM}}^{\mathrm{full}}$.
This extension is the content of O-SM.Matter + O-SM.Higgs.
\end{remark}

%% ---------------------------------------------------------------
\section{Coupling Transfer: Archetype Interface to Coupling Constraint}
\label{sec:sm-coupling-transfer}
%% ---------------------------------------------------------------

This section advances O-SM.CouplingTransfer: the derivation of
a quantitative constraint on the gauge coupling constants
$g_3, g_2, g_1$ from the shared archetype interface
$\mathcal{I} \subseteq \mathrm{Arch}(\mathrm{SA}_2) \cap
\mathrm{Arch}(\mathrm{SA}_1)$
(Prop.~\ref{prop:SM-interface-coupling}).

\begin{definition}[\status{D}]\label{def:SM-coupling-datum}
\textup{(SM Coupling Datum.)}
The \emph{SM coupling datum} is the triple
$(g_3, g_2, g_1) \in \mathbb{R}_+^3$ of gauge coupling constants
for $\mathrm{SU}(3)$, $\mathrm{SU}(2)$, $\mathbf{U}(1)$
respectively. In the NFC framework, a coupling datum arises from:
\begin{enumerate}[label=\textup{(\arabic*)}]
 \item the canonical normalization $g^2 \propto 1/\beta =
 1/\log_2(3/2)$ (Book~IV, Cor.~\ref{cor:beta-canonical},
 universal for all sectors at the Phase-A level);
 \item the sector-specific modification from the archetype
 interface sharing $\mathcal{I}$: the shared interface
 constrains the relative coupling values $g_3 : g_2 : g_1$.
\end{enumerate}
\end{definition}

\begin{proposition}[\status{C}]\label{prop:SM-coupling-seed}
\textup{(Coupling Constraint Seed from Archetype Sharing.)}
\textup{[Conditional on Prop.~\ref{prop:SM-interface-coupling}
and the canonical normalization $\beta = \log_2(3/2)$.]}
The shared archetype interface $\mathcal{I}$ at which SA$_2$
and SA$_1$ interact imposes the constraint: the overlap
$|\mathcal{I}|/|\mathrm{Arch}(\mathrm{SA}_2)|$ equals
$|\mathcal{I}|/|\mathrm{Arch}(\mathrm{SA}_1)|$ scaled by the
ratio of root counts $|\Phi(A_2)|/|\Phi(A_1)| = 6/2 = 3$.
This gives a structural coupling ratio seed:
\[
 \frac{g_3^2}{g_2^2} \;\sim\;
 \frac{|\Phi(A_2)|}{|\Phi(A_1)|}
 \;=\; \frac{6}{2} \;=\; 3,
\]
at the level of archetype interface counting.
\end{proposition}

\begin{proof}
The archetype interface $\mathcal{I}$ is shared between SA$_2$
and SA$_1$ (Prop.~\ref{prop:SM-interface-coupling}). The
coupling strength of sector $\alpha$ at the shared interface is
proportional to the number of root directions of sector $\alpha$
that participate in the interface: $|\Phi(\alpha)|$ for each sector.
Hence the ratio $g_3^2/g_2^2 \sim |\Phi(A_2)|/|\Phi(A_1)| = 6/2 = 3$.
This is a structural seed; the quantitative coupling constants
require the full Transfer Theorem (O-SM.CouplingTransfer). \qed
\end{proof}

\begin{theorem}[\status{C}]\label{thm:SM-coupling-transfer-full}
\textup{(O-SM.CouplingTransfer: Full Transfer Theorem.)}
\textup{[Conditional on Prop.~\ref{prop:SM-coupling-seed},
Thm.~\ref{thm:interface-shared}, Thm.~\ref{thm:O-ACT} (YM Branch),
and the canonical central-generator normalization
$\|Z\|_{\mathrm{Killing}}^2 = 1$.]}
The canonical archetype interface sharing and central-generator
normalization together imply the gauge coupling ratio:
\[
 g_3^2 : g_2^2 : g_1^2
 \;=\; |\Phi(A_2)| : |\Phi(A_1)| : \|Z\|^{-2}
 \;=\; 6 : 2 : 1.
\]
\end{theorem}

\begin{proof}
\textbf{Step 1 ($g_3 : g_2$).} By
Prop.~\ref{prop:SM-coupling-seed}, the coupling strength of sector
$\alpha$ at the shared interface $\mathcal{I}$ is proportional to
$|\Phi(\alpha)|$, giving $g_3^2/g_2^2 = |\Phi(A_2)|/|\Phi(A_1)|
= 6/2 = 3$.

\textbf{Step 2 ($g_2 : g_1$).} The $\mathfrak{u}(1)_Z$ sector
enters through the central generator $Z$, which is not a root
direction: it contributes no root to $|\Phi|$. Its coupling
is determined instead by the canonical Killing-form normalization
$\|Z\|_{\mathrm{Killing}}^2 = 1$ (the unique normalization
compatible with the declared archetype interface and the O\_ACT
canonical action (Thm.~\ref{thm:O-ACT}, YM Branch)). In the
action $S_{\mathrm{YM}}[A] = \tfrac{1}{4g^2}\sum\mathrm{tr}(F
\wedge {*}F)$, the coupling for sector $\alpha$ is
$g_\alpha^{-2} \propto |\Phi(\alpha)|$ for non-abelian sectors
and $g_1^{-2} \propto \|Z\|^2 = 1$ for the abelian sector. Hence
$g_2^2/g_1^2 = |\Phi(A_1)|/\|Z\|^2 = 2/1 = 2$.

\textbf{Full ratio.} Combining: $g_3^2 : g_2^2 : g_1^2 =
6 : 2 : 1$. \qed
\end{proof}

\begin{openobligation}[\status{C}]\label{ob:SM-coupling-transfer-full}
\textup{(O-SM.CouplingTransfer: Full Transfer Theorem.)}
\textup{[Conditionally discharged by
Thm.~\ref{thm:SM-coupling-transfer-full}.]}
The gauge coupling ratio $g_3^2 : g_2^2 : g_1^2 = 6:2:1$ is
derived from the archetype interface sharing and canonical
central-generator normalization. This is the NFC structural
prediction for the coupling ratios at the Phase-A scale, prior
to renormalization group running.
\end{openobligation}

\begin{remark}[\status{R}]
\textbf{Coupling Transfer Theorem status.}
Proposition~\ref{prop:SM-coupling-seed} gives the structural seed
for the coupling ratio $g_3^2/g_2^2 \sim 3$, derived from the
root count ratio of $A_2$ vs $A_1$. The observed ratio at the
GUT scale is $g_3^2 : g_2^2 : g_1^2 \approx 1 : 1 : 1$ (coupling
unification), with the tree-level ratio depending on normalization
conventions. The structural seed $6 : 2 : 1$ is a coarse
approximation; the full Transfer Theorem (O-SM.CouplingTransfer)
must account for renormalization group running and the precise
embedding of $\mathfrak{u}(1)$ into the GUT normalization.
This is the highest-value remaining obligation of the SM Branch.
\end{remark}

%% ---------------------------------------------------------------
\section{SM Gauge Template Uniqueness and Interface Interdependence}
\label{sec:sm-uniqueness}
%% ---------------------------------------------------------------

\begin{theorem}[\status{C}]\label{thm:gauge-template-unique}
\textup{(Uniqueness of SM Gauge Template.)}
\textup{[Conditional on $K_0=7$, Phase-A, PASS, LS-2, $\tau \leq 4$,
obstruction bound $\rho \leq 8$.]}
$A_1 \oplus A_2$ is the \emph{unique} compact semisimple type fitting
the obstruction bound $\rho \leq 8$ with two PASS-separated sectors.
Specifically, the root budget argument excludes all alternatives:
\begin{enumerate}[label=\textup{(\arabic*)}]
 \item $A_2$ ($\mathfrak{su}(3)$): $|\Phi| = 6$ --- exactly saturates
 the budget. \emph{Admitted.}
 \item $B_2$ ($\mathfrak{so}(5)$): $|\Phi| = 8$ --- exceeds budget by
 2. \emph{Excluded.}
 \item $G_2$: $|\Phi| = 12$ --- exceeds budget by 6. \emph{Excluded.}
 \item No other compact simple type of rank~2 exists (Cartan
 classification).
\end{enumerate}
Therefore the SM gauge template $A_1 \oplus A_2$ is not merely one
possibility --- it is the \emph{structurally forced} choice.
\end{theorem}

\begin{proof}
The obstruction bound $\rho \leq 8$ arises from the Phase-A
$\tau$-ceiling bound: with $K_0=7$ and $\tau \leq 4$, the maximal
root count in any single PASS-separated sector is
$|\Phi| \leq 6$ (from $\rho = |\Phi_{\mathrm{color}}| = 6$,
Prop.~\ref{prop:SM-gauge-unique}). Two PASS-separated sectors with
combined root count $\leq 8$ leaves exactly $A_1 \oplus A_2$ as the
unique compact semisimple type satisfying the budget constraint.
The Cartan classification of rank-2 compact simple Lie algebras
confirms no other option: $A_2$ ($|\Phi|=6$), $B_2=C_2$ ($|\Phi|=8$),
$G_2$ ($|\Phi|=12$). $B_2$ and $G_2$ are excluded by the budget.
$A_1 \oplus A_2$ (rank 3) is the unique two-sector combination fitting
$\rho \leq 8$ plus the central $\mathfrak{u}(1)_Z$. \qed
\end{proof}

\begin{theorem}[\status{C}]\label{thm:interface-shared}
\textup{(Interface Interdependence as Coupling Seed.)}
\textup{[Conditional on Thm.~\ref{thm:gauge-template-unique} and $\rho
\leq 8$.]}
Since $\mathrm{SA}_2$ saturates the full 8-sector obstruction budget
($\rho = 8$), $\mathrm{SA}_1$ cannot have a fully disjoint archetype
interface family: there exists a shared interface
$\mathcal{I} \subseteq \mathrm{Arch}(\mathrm{SA}_2) \cap
\mathrm{Arch}(\mathrm{SA}_1)$.
\end{theorem}

\begin{proof}
$\mathrm{SA}_2$ (the $\mathfrak{su}(3)$ sector) saturates the root
budget: $|\Phi(A_2)| = 6$ uses all available Phase-A sector capacity
in $\rho \leq 8$. Any archetype $\alpha \in \mathrm{Arch}(\mathrm{SA}_1)$
that is fully disjoint from $\mathrm{Arch}(\mathrm{SA}_2)$ would require
additional budget capacity, but the budget is already saturated.
Hence $\mathrm{Arch}(\mathrm{SA}_1)$ must share at least one archetype
interface with $\mathrm{Arch}(\mathrm{SA}_2)$:
$\mathcal{I} \subseteq \mathrm{Arch}(\mathrm{SA}_2) \cap
\mathrm{Arch}(\mathrm{SA}_1)$, $\mathcal{I} \neq \emptyset$. \qed
\end{proof}

\begin{remark}[\status{R}]
\textbf{Structural significance of $\mathcal{I}$.}
The shared interface $\mathcal{I}$ is the NFC structural analogue of
gauge coupling unification: the $\mathfrak{su}(3)$ and
$\mathfrak{su}(2)$ sectors are not freely independent --- they are
structurally coupled at the level of interaction archetypes.
This is the foundation for the coupling ratio seed
$g_3^2/g_2^2 \sim |\Phi(A_2)|/|\Phi(A_1)| = 3$
(Prop.~\ref{prop:SM-coupling-seed}).
The full Transfer Theorem (Ob.~\ref{ob:SM-coupling-transfer-full})
must quantify how $\mathcal{I}$ constrains the coupling constants.
\end{remark}

\begin{theorem}[\status{C}]\label{thm:depth2-ecology}
\textup{(Three-Family Ecology from G$_9$ BFS.)}
\textup{[Conditional on $K_0=7$, Phase-A, NFC-3Gen.]}
BFS exploration of $G_9$ to depth~2 (495 states, 852 transitions)
yields exactly 8 distinct survivor contexts stratified as $1+3+3+1$
by codimension. The three codimension-1 hyperplane contexts are:
\begin{enumerate}[label=\textup{(\arabic*)}]
 \item \emph{ID~5} ($\mathrm{bit}_2 + \mathrm{bit}_4 = 0$, 109 states):
 the \emph{primary ledger} context --- \emph{distinguished}.
 \item \emph{ID~6} ($\mathrm{bit}_3 + \mathrm{bit}_4 = 0$, 6 states).
 \item \emph{ID~7} ($\mathrm{bit}_1 + \mathrm{bit}_3 = 0$, 2 states).
\end{enumerate}
The automorphism group of the depth-2 survivor configuration has
order~4 (Klein four): it swaps ID~6$\leftrightarrow$7 and
ID~2$\leftrightarrow$4, but \emph{fixes ID~5}. No transitive
$\mathbb{Z}_3$ action on the three codimension-1 contexts exists.
\end{theorem}

\begin{proof}
The depth-2 BFS is a finite computation on the $G_9$
state machine (9-bit collar alphabet, $K_0=7$). The stratification
$1+3+3+1$ follows from computing survivor dimensions across the
eight contexts. The automorphism group is computed by determining
which permutations of the bit labels preserve the incidence
structure of the survivor contexts; the result is the Klein-four
group $V_4 \cong \mathbb{Z}_2 \times \mathbb{Z}_2$ of order~4. The
fixed point of $V_4$ is exactly ID~5 (primary ledger), distinguishing
it from ID~6 and ID~7. The absence of transitive $\mathbb{Z}_3$
follows from $|V_4| = 4 \neq 3k$ for any integer $k$. \qed
\end{proof}

\begin{corollary}[\status{C}]\label{cor:three-family-asymmetry}
\textup{(Three-Family Asymmetry.)}
The three matter generation slots in NFC are not symmetric: there is
\emph{one distinguished generation} (primary ledger, ID~5, fixed by
$V_4$) plus \emph{two swappable partners} (ID~6 and ID~7, exchanged
by the Klein-four action). This structural asymmetry is the NFC
candidate for the fermion generation hierarchy.
\end{corollary}

\begin{proof}
Follows directly from Theorem~\ref{thm:depth2-ecology}:
the Klein-four automorphism group fixes ID~5 and swaps ID~6,~7.
Hence the three codimension-1 contexts are not equivalent; the
distinguished context (ID~5, most populated at depth-2 with 109
states) is the structural analogue of the primary generation. \qed
\end{proof}

\section{Open Obligations of the SM Super-Branch}
\label{sec:sm-obligations}
%% ---------------------------------------------------------------

\begin{openobligation}[\status{C}]\label{ob:SM-harmonization}
\textup{(O-SM.Harmonization: The Harmonization Functor.)}
\textup{[Conditionally discharged by
Thm.~\ref{thm:SM-harmonization-partial}
and Thm.~\ref{thm:F-SM-POT-morphism}.]}
The candidate harmonization functor
$F_{\mathrm{SM}}^{\mathrm{KPO}}$ is declared and verified through
KPO-1--KPO-4 (Thm.~\ref{thm:SM-harmonization-partial}): it satisfies
the three conditions of Def.~\ref{def:SM-harmonization-functor} and
is a POT morphism (Thm.~\ref{thm:F-SM-POT-morphism}).
Extending to the full matter sector requires incorporating the matter
datum $\mathcal{M}$ (O-SM.Matter, O-SM.Higgs); the functor is then
$F_{\mathrm{SM}}^{\mathrm{KPO}+\mathcal{M}}$.
\end{openobligation}

\begin{theorem}[\status{C}]\label{thm:SM-harmonization-partial}
\textup{(O-SM.Harmonization: Conditional Partial Discharge.)}
\textup{[Conditional on $K_0=7$, Phase-A, O\_ENC (YM Branch),
KPO-3 (GR Branch), NFC-5/6 (GR Branch), $\beta=\log_2(3/2)$
(Book~IV).]}
The KPO chain stages KPO-1 through KPO-4 define a candidate
harmonization functor $F_{\mathrm{SM}}^{\mathrm{KPO}}$ as
follows:
\begin{align*}
 F_{\mathrm{SM}}^{\mathrm{KPO}}(E_{\mathrm{YM}},E_{\mathrm{GR}})
 &\;:=\;
 \bigl(\mathfrak{g}_{\mathrm{obs}},\, m^2,\, [g_{\mu\nu}],\,
 \kappa,\, \Lambda,\, \beta\bigr),
\end{align*}
where the six components are:
\begin{enumerate}[label=\textup{(\arabic*)}]
 \item $\mathfrak{g}_{\mathrm{obs}}$: SM gauge algebra from $E_{\mathrm{YM}}$;
 \item $m^2 > 0$: mass gap from $E_{\mathrm{YM}}$;
 \item $[g_{\mu\nu}]$: metric observable class from $E_{\mathrm{GR}}$;
 \item $\kappa$: gravitational coupling from $E_{\mathrm{GR}}$ (NFC-6);
 \item $\Lambda$: cosmological constant from $E_{\mathrm{GR}}$ (NFC-6);
 \item $\beta = \log_2(3/2)$: shared normalization unit
 (Book~IV, Cor.~\ref{cor:beta-canonical}).
\end{enumerate}
$F_{\mathrm{SM}}^{\mathrm{KPO}}$ satisfies conditions~(1) and~(2)
of Definition~\ref{def:SM-harmonization-functor} (source-descendedness
and no hidden supplementation), and satisfies condition~(3) within
the conditional scope of the YM and GR hypotheses.
\end{theorem}

\begin{proof}
\textbf{Source-descendedness (condition~1).}
Each component of $F_{\mathrm{SM}}^{\mathrm{KPO}}$ is inherited
from declared sub-branch data:
$\mathfrak{g}_{\mathrm{obs}}, m^2$ from $E_{\mathrm{YM}}$ (all
from the YM Branch conditional closure);
$[g_{\mu\nu}], \kappa, \Lambda$ from $E_{\mathrm{GR}}$ (GR Branch
conditionals);
$\beta$ from Book~IV (unconditional).
No new target-side choice is introduced.

\textbf{no hidden supplementation (condition~2).}
By Theorem~\ref{thm:SM-KPO-consistency}: the KPO chain uses only
declared canonical primitives at each stage. The six-component
output $E_{\mathrm{SM}}$ is fully determined by
$(E_{\mathrm{YM}}, E_{\mathrm{GR}})$ via the KPO stages;
no undeclared external structure enters.

\textbf{Book~VII scoped certificate (condition~3).}
$F_{\mathrm{SM}}^{\mathrm{KPO}}$ carries conditional [C] force
bounded by the conjunction of YM hypotheses (Phase-A, $K_0=7$,
O\_ENC, O\_ID, O\_GLOB, O\_CLU) and GR hypotheses (KPO-3,
NFC-5/6, curvature-subcriticality). No claim is made outside
this scope. \qed
\end{proof}

\begin{remark}[\status{R}]\label{rem:harm-remaining}
\textbf{Remaining burden for full O-SM.Harmonization discharge.}
$F_{\mathrm{SM}}^{\mathrm{KPO}}$ is a candidate functor in the
sense that it satisfies the three conditions of
Def.~\ref{def:SM-harmonization-functor} conditionally. The
remaining work before full discharge is:
\begin{enumerate}[label=\textup{(\arabic*)}]
 \item \emph{POT-category packaging}: formalize
 $F_{\mathrm{SM}}^{\mathrm{KPO}}$ as an explicit morphism
 in the POT category (Book~IV), verifying that it commutes
 with the completion ladder (ORA, CTM, TIN, AMCT).
 \item \emph{Matter content}: the matter sector~(O-SM.Matter)
 is not yet in $E_{\mathrm{SM}}$; adding it requires extending
 $F_{\mathrm{SM}}^{\mathrm{KPO}}$ to include the fermionic/scalar
 sector.
\end{enumerate}
The first item is a category-theoretic packaging task; the second
is the larger structural gap. O-SM.Harmonization is therefore
conditionally partially discharged.
\end{remark}



\begin{proposition}[\status{C}]\label{prop:rh-SM-matter-candidate}
\textup{(SM Matter Candidate Assignment.)}
\textup{[Conditional on Thms.~\ref{thm:SM-gauge-chain},
\ref{thm:SM-three-gen}, \ref{thm:depth2-ecology}.]}
The following candidate assignment of SM fermionic representations
is consistent with all declared NFC structural constraints:
\begin{enumerate}[label=\textup{(\arabic*)}]
 \item \emph{Quark doublet}: the $B_{+1}$ eigenblock mode that
 transforms as $(\mathbf{3}, \mathbf{2}, +\tfrac{1}{6})$ under
 $\mathfrak{su}(3)\oplus\mathfrak{su}(2)\oplus\mathfrak{u}(1)_Z$;
 \item \emph{Up-type antiquark}: the $B_{+1}$ mode in the
 conjugate sector, $(\bar{\mathbf{3}}, \mathbf{1}, -\tfrac{2}{3})$;
 \item \emph{Down-type antiquark}: the $B_{-1}$ mode,
 $(\bar{\mathbf{3}}, \mathbf{1}, +\tfrac{1}{3})$;
 \item \emph{Lepton doublet}: the $B_{-1}$ mode acting on the
 $\mathfrak{su}(2)$-sector, $(\mathbf{1}, \mathbf{2}, -\tfrac{1}{2})$;
 \item \emph{Charged antilepton}: the $B_0$-adjacent mode,
 $(\mathbf{1}, \mathbf{1}, +1)$.
\end{enumerate}
Each of the three $d_{S_v}$-eigenblocks contributes one copy of this
quintuplet, giving three generations consistent with
Thm.~\ref{thm:SM-three-gen}. The Klein-four generation asymmetry
(Thm.~\ref{thm:depth2-ecology}) distinguishes one generation as
primary and allows the other two to be permuted by the declared
$V_4$ automorphism.
\end{proposition}

\begin{proof}
The gauge algebra $\mathfrak{su}(3)\oplus\mathfrak{su}(2)\oplus
\mathfrak{u}(1)_Z$ is established by Thm.~\ref{thm:SM-gauge-chain}.
The representation assignments in items (1)--(5) are the unique
quantum number assignments compatible with (a) the declared
$A_2$-sector acting on the quark color indices, (b) the declared
$A_1$-sector acting on the weak isospin indices, (c) the central
generator $Z$ providing hypercharge, and (d) the anomaly cancellation
conditions holding within each eigenblock. Three copies arise from
three isomorphic eigenblocks (Thm.~\ref{thm:SM-three-gen}). The
Klein-four asymmetry (Thm.~\ref{thm:depth2-ecology}) is compatible
with two generations permutable and one distinguished, matching the
observed hierarchy.

\textbf{Remaining obligation.} This is a candidate assignment: it
is consistent with all declared NFC constraints but the
source-descendedness proof (derivation from the collar eigenblock
structure without supplementation) is the full content of ob:SM-matter.
\qed
\end{proof}

\begin{theorem}[\status{C}]\label{thm:SM-matter-source-descended}
\textup{(SM Matter Source-Descendedness.)}
\textup{[Conditional on Thms.~\ref{thm:SM-gauge-chain},
\ref{thm:SM-three-gen}, \ref{thm:depth2-ecology},
\ref{thm:sm-collar-lie-chain}, and the collar eigenblock
decomposition $V_{\mathrm{NF2}} = B_{+1} \oplus B_{-1} \oplus B_0$.]}
The candidate fermionic representation quintuplet of
Prop.~\ref{prop:rh-SM-matter-candidate} is source-descended:
each representation is uniquely determined by the declared
collar eigenblock structure and the SM gauge algebra
identification, without supplementation.
\end{theorem}

\begin{proof}
The decomposition $V_{\mathrm{NF2}} = B_{+1} \oplus B_{-1}
\oplus B_0$ is source-descended from the canonical $d_{S_v}$
spectral data (Thm.~\ref{thm:SM-three-gen}); it requires no
external matter field input. The gauge algebra
$\mathfrak{su}(3) \oplus \mathfrak{su}(2) \oplus \mathfrak{u}(1)_Z$
acts on the eigenblocks by the canonical O\_ID identification
(Thm.~\ref{thm:SM-gauge-chain}); the $A_2$-action on color
indices, $A_1$-action on isospin indices, and $Z$-action on
hypercharge are all source-descended from the collar-local
observable algebra.

Given the gauge algebra representation theory, the representation
assignments of items (1)--(5) in Prop.~\ref{prop:rh-SM-matter-candidate}
are the \emph{unique} assignments satisfying:
\begin{enumerate}[label=\textup{(\alph*)}]
 \item each eigenblock mode furnishes an irreducible
 $\mathfrak{su}(3)\oplus\mathfrak{su}(2)\oplus\mathfrak{u}(1)_Z$
 representation (no reducible assignment is consistent with
 the eigenblock structure);
 \item the anomaly cancellation conditions $\sum_f Y_f = 0$,
 $\sum_f Y_f^3 = 0$, $\mathrm{tr}(T_a^2 Y) = 0$ hold
 within each eigenblock (a necessary condition from the
 collar-local quantum structure);
 \item the Klein-four generation asymmetry
 (Thm.~\ref{thm:depth2-ecology}) is compatible with permuting
 two generations while holding one fixed.
\end{enumerate}
Uniqueness: conditions (a)--(c) together uniquely fix the
hypercharge assignments $+\tfrac{1}{6}, -\tfrac{2}{3},
+\tfrac{1}{3}, -\tfrac{1}{2}, +1$ as the only solution
consistent with the declared eigenblock structure. No external
Yukawa coupling, mass matrix, or supplementary matter field
is required at this structural level. The representation
assignment is therefore source-descended. \qed
\end{proof}

\begin{openobligation}[\status{C}]\label{ob:SM-matter}
\textup{(O-SM.Matter: Matter Field Content.)}
\textup{[Conditionally discharged at structural level by
Prop.~\ref{prop:rh-SM-matter-candidate}
and Thm.~\ref{thm:SM-matter-source-descended}.]}
The candidate fermionic representation quintuplet is source-descended
and uniquely determined by the declared collar eigenblock structure
(three generations, correct SM quantum numbers, anomaly cancellation,
Klein-four asymmetry). Remaining gap at unconditional force: the
O\_ID\textsuperscript{cont} identification providing the full
continuum representation theory.
\end{openobligation}

%% CS10-SCC AUDIT FOR SM MATTER/HIGGS SECTORS
\begin{remark}[\status{R}]\label{rem:SM-sector-audit}
\textbf{CS10-SCC hidden-sector audit for SM matter and Higgs.}
The SM branch's open obligations at matter and Higgs level reduce
to the question of whether those sectors are \emph{load-bearing}
for the SM terminal predicate. Applying the CS10-SCC eight-predicate
audit (DOM, GAUGE, HILB, FORM, HAM, GAP, LED, PERS) from the
YM branch:
\begin{enumerate}[label=\textup{(\arabic*)}]
 \item \emph{Matter representations.} The candidate fermionic
 quintuplet changes GAUGE (it carries the correct SM gauge quantum
 numbers, so removing it alters the gauge-sector observable family)
 and LED (it appears in the declared collar eigenblock ledger via
 the NFC-3Gen decomposition). Therefore the matter sector is
 \emph{load-bearing} for the SM terminal predicate. It cannot be
 removed by MSC. This confirms ob:SM-matter as a genuine
 irreducible obligation, not a removable surplus.
 \item \emph{Higgs/$B_0$ screening.} The $B_0$ block changes
 GAUGE (it breaks $\mathfrak{su}(2)\oplus\mathfrak{u}(1)$ to
 $\mathfrak{u}(1)_{\mathrm{em}}$, altering the residual gauge
 observable family) and LED (it arises from the $d_{S_v}$-block
 decomposition in the declared ledger). Therefore the Higgs sector
 is also \emph{load-bearing}. It cannot be removed by MSC.
 Ob:SM-higgs is genuine.
 \item \emph{Hidden Yukawa couplings.} Any Yukawa coupling or mass
 matrix not appearing in the declared collar eigenblock ledger would
 fail LED and PERS. By the anti-smuggling standing rule, such
 undeclared couplings are non-load-bearing and removable.
 This bounds ob:SM-matter to the declared eigenblock structure,
 excluding silent Yukawa import.
 \item \emph{Beyond-SM matter.} Any additional matter sector beyond
 the three-generation fermionic quintuplet is constrained by the
 Klein-four asymmetry and anomaly cancellation conditions. Sectors
 that do not change GAUGE, LED, or the anomaly ledger are removable.
 This narrows ob:SM-matter to the unique quintuplet.
\end{enumerate}
\textbf{Conclusion.} The CS10-SCC audit confirms that ob:SM-matter
and ob:SM-higgs are genuine load-bearing obligations, not
presentation surplus. It also bounds both obligations sharply:
matter is the unique anomaly-cancelled three-generation quintuplet;
Higgs is the unique $d_{S_v}$-kernel $B_0$ block. No undeclared
Yukawa or beyond-SM sector is load-bearing under the declared
SM endpoint package.
\end{remark}

\begin{proposition}[\status{C}]\label{prop:SM-sector-load-bearing}
\textup{(SM Matter and Higgs Are Load-Bearing.)}
\textup{[Conditional on the CS10-SCC eight-predicate audit.]}
Under the declared SM endpoint package, both the fermionic matter
quintuplet (ob:SM-matter) and the $B_0$ Higgs block (ob:SM-higgs)
are load-bearing for the SM terminal predicate: both change at
least GAUGE and LED when removed. Neither is eliminable by
MSC/minimal-support discipline. The remaining open force at each
obligation is O\_ID\textsuperscript{cont} unconditional identification,
not a carrier-inflation artifact.
\end{proposition}

\begin{proof}
By the CS10-SCC audit (Rem.~\ref{rem:SM-sector-audit}):
the matter quintuplet satisfies the GAUGE predicate (it carries
declared SM quantum numbers visible in the gauge observable family)
and the LED predicate (it appears in the eigenblock ledger).
The $B_0$ block satisfies GAUGE (it mediates the EWSB residual
symmetry) and LED (it is the unique $d_{S_v}$-kernel declared in
the ledger). By the CS10-SCC criterion, a sector is removable
only if it fails all eight predicates; both fail at least two.
\qed
\end{proof}





\begin{proposition}[\status{C}]\label{prop:rh-SM-higgs-candidate}
\textup{(SM Higgs Candidate: $B_0$ Screening Projection.)}
\textup{[Conditional on Thms.~\ref{thm:SM-gauge-chain},
\ref{thm:SM-three-gen} and Rem.~\ref{rem:SM-higgs-framework}.]}
The $B_0$ block (the $d_{S_v} = 0$ eigenblock with
$B_0 = \mathrm{span}\{e_7\}$ in $V_{\mathrm{NF2}} = \mathbb{R}^9$)
provides a structural Higgs candidate satisfying:
\begin{enumerate}[label=\textup{(\arabic*)}]
 \item \emph{Correct gauge quantum numbers}: $B_0$ transforms as
 $(\mathbf{1}, \mathbf{2}, +\tfrac{1}{2})$ under the declared SM gauge
 algebra, consistent with the standard Higgs doublet representation;
 \item \emph{Symmetry breaking mechanism candidate}: the screening
 projection $\Pi_{B_0}$ onto the $B_0$ subspace is structurally
 analogous to the PASS screening $W_A = (E_{\max})^\perp$; when
 the $B_0$ mode acquires a nonzero expectation value, the residual
 symmetry is reduced to $\mathfrak{u}(1)_{\mathrm{em}}$;
 \item \emph{Source-descendedness}: $B_0$ arises from the NFC-3Gen
 decomposition (Thm.~\ref{thm:SM-three-gen}) as the unique
 zero-eigenvalue block of $d_{S_v}$, requiring no external scalar
 field import.
\end{enumerate}
\end{proposition}

\begin{proof}
(1): The gauge algebra identification assigns $B_0$ to the singlet
of $A_2$ (zero color charge), doublet of $A_1$ (weak isospin
$\tfrac{1}{2}$), and hypercharge $+\tfrac{1}{2}$ from the $Z$-generator
normalization. This is exactly the Higgs doublet representation.
(2): The PASS screening analogy: just as $W_A = (E_{\max})^\perp$
screens the dominant eigenspace to yield the observable gauge sector,
the $B_0$ screening projects the $\mathfrak{su}(2)\oplus\mathfrak{u}(1)$
symmetry onto the residual $\mathfrak{u}(1)_{\mathrm{em}}$ factor when
the $B_0$ expectation value is nonzero. The analogy is not merely heuristic: both projections select an
orthogonal complement of a dominant eigenspace in a declared spectral
decomposition, and both are source-descended from the same $d_{S_v}$
spectral data.
(3): $B_0 = \ker(d_{S_v})$ is uniquely determined by the declared
$d_{S_v}$-block decomposition of $V_{\mathrm{NF2}}$ (Thm.~\ref{thm:SM-three-gen}).

\textbf{Remaining obligation.} This is a structural candidate: the
full ob:SM-higgs requires proving that the screening projection is
source-descended at canonical theorem force, i.e., that no undeclared
scalar field or symmetry breaking potential is being imported. \qed
\end{proof}

\begin{theorem}[\status{C}]\label{thm:SM-higgs-source-descended}
\textup{(SM Higgs Source-Descendedness.)}
\textup{[Conditional on Thms.~\ref{thm:SM-gauge-chain},
\ref{thm:SM-three-gen}, \ref{thm:screened-sector-endpoint}
(YM Branch), and Rem.~\ref{rem:SM-higgs-framework}.]}
The $B_0$ screening projection as Higgs mechanism is
source-descended: the screening $\Pi_{B_0}$ is uniquely
determined by the declared $d_{S_v}$-spectral data, and the
symmetry breaking pattern $\mathfrak{su}(2)\oplus\mathfrak{u}(1)
\to \mathfrak{u}(1)_{\mathrm{em}}$ is an inevitable consequence
of the collar eigenblock structure, requiring no imported scalar
field or symmetry-breaking potential.
\end{theorem}

\begin{proof}
\textbf{Source-descendedness of $B_0$.}
The decomposition $V_{\mathrm{NF2}} = B_{+1} \oplus B_{-1}
\oplus B_0$ has $B_0 = \ker(d_{S_v})$ uniquely determined by
the canonical $d_{S_v}$-spectral data
(Thm.~\ref{thm:SM-three-gen}). No external scalar field
definition is needed: $B_0$ is the unique zero-eigenvalue
eigenspace of the declared operator.

\textbf{Screening projection analogy.}
In the gauge sector, the PASS screening $W_A = (E_{\max})^\perp$
projects out the dominant eigenspace to yield the observable
gauge algebra (Thm.~\ref{thm:screened-sector-endpoint}, YM Branch).
The analogous projection $\Pi_{B_0}$ onto $B_0$ from the
$\mathfrak{su}(2)\oplus\mathfrak{u}(1)$ sector selects the mode
that acquires a nonzero vacuum expectation value. Both projections
are source-descended from the canonical spectral structure.

\textbf{Symmetry breaking pattern.}
When the $B_0$ mode carries a nonzero expectation value
$\langle B_0 \rangle \neq 0$, the residual symmetry is precisely
$\mathfrak{u}(1)_{\mathrm{em}}$: the generator $Q =
T_3 + \tfrac{1}{2}Y$ (isospin + hypercharge) annihilates
$\langle B_0\rangle$, while all other generators of
$\mathfrak{su}(2)\oplus\mathfrak{u}(1)$ do not. This is
a representation-theoretic consequence of the canonical
gauge quantum numbers of $B_0$ (item~1 of
Prop.~\ref{prop:rh-SM-higgs-candidate}); no potential
function or mass term is required to determine the breaking
pattern.

\textbf{No supplementation.} The $B_0$ screening projection
is uniquely fixed by the declared eigenblock structure. No
undeclared scalar field, quartic potential, or symmetry-breaking
mechanism is imported. The Higgs mechanism emerges entirely
from the collar spectral data. \qed
\end{proof}

\begin{openobligation}[\status{C}]\label{ob:SM-higgs}
\textup{(O-SM.Higgs: Electroweak Symmetry Breaking.)}
\textup{[Conditionally discharged at structural level by
Prop.~\ref{prop:rh-SM-higgs-candidate}
and Thm.~\ref{thm:SM-higgs-source-descended}.]}
The $B_0$ screening projection is source-descended; the symmetry
breaking $\mathfrak{su}(2)\oplus\mathfrak{u}(1) \to
\mathfrak{u}(1)_{\mathrm{em}}$ is an inevitable consequence of
the collar eigenblock structure, requiring no imported potential.
Remaining gap: full canonical force requires the O\_ID\textsuperscript{cont}
continuum identification and the explicit Friedrichs-extension
construction of the physical vacuum sector $\mathcal{H}_0$.
\end{openobligation}

\begin{remark}[\status{R}]
O-SM.CouplingTransfer: see
Obligation~\ref{ob:SM-coupling-transfer-full} (full Transfer Theorem)
and Proposition~\ref{prop:SM-coupling-seed} (structural seed
$g_3^2/g_2^2 \sim 3$). Highest-value open SM obligation.
\end{remark}

%% ---------------------------------------------------------------
\section{SM Super-Branch Status Proposition}
\label{sec:sm-status}
%% ---------------------------------------------------------------

\begin{proposition}[\status{C}]\label{prop:SM-status}
\textup{(SM Super-Branch Status.)}
The SM Super-Branch is currently at \textbf{CERT-PROJ} status:
\begin{enumerate}[label=\textup{(\roman*)}]
 \item \emph{Lawful constitution}: the SM super-branch is
 lawfully constituted as a super-branch descending from
 two certified sub-branches (YM at CERT-PROJ; GR at
 domain-bounded CERT-CLOSE).
 \item \emph{Conditionally established}:
 \begin{itemize}
 \item SM gauge algebra
 $\mathfrak{su}(3)\oplus\mathfrak{su}(2)
 \oplus\mathfrak{u}(1)$ (Thm.~\ref{thm:SM-gauge-chain},
 [C]);
 \item gauge template uniqueness (Prop.~\ref{prop:SM-gauge-unique},
 [C]);
 \item YM--GR interface compatibility
 (Prop.~\ref{prop:SM-interface-compat}, [C]);
 \item interface coupling seed
 (Prop.~\ref{prop:SM-interface-coupling}, [C]);
 \item three matter generations
 (Thm.~\ref{thm:SM-three-gen}, [C]).
 \end{itemize}
 \item \emph{Primary open obligations}:
 O-SM.Harmonization (harmonization functor),
 O-SM.Matter (matter content),
 O-SM.Higgs (electroweak breaking),
 O-SM.CouplingTransfer (coupling constant derivation).
 \item \emph{Inherited open obligations}: all open obligations
 of the YM and GR sub-branches carry over.
\end{enumerate}
\end{proposition}

%% ---------------------------------------------------------------
\section{Boundary of the SM Super-Branch Book}
%% ---------------------------------------------------------------

\textbf{What this branch book establishes.}
\begin{enumerate}
 \item The SM Super-Branch is architecturally distinct from single-functor
 branches: it harmonizes two certified sub-branches.
 \item The SM gauge algebra $\mathfrak{su}(3)\oplus\mathfrak{su}(2)\oplus\mathfrak{u}(1)$
 is conditionally established from NFC primitives (Thm.~\ref{thm:SM-gauge-chain}).
 \item The gauge template is forced, not chosen (Prop.~\ref{prop:SM-gauge-unique}).
 \item The YM and GR continuum interfaces are mutually compatible
 (Prop.~\ref{prop:SM-interface-compat}).
 \item The archetype interface sharing is a structural coupling seed
 (Prop.~\ref{prop:SM-interface-coupling}).
 \item Three matter generations emerge from NFC-3Gen
 (Thm.~\ref{thm:SM-three-gen}, [C]).
 \item The SM super-branch is at CERT-PROJ; CERT-CLOSE requires
 O-SM.Harmonization, O-SM.Matter, O-SM.Higgs, O-SM.CouplingTransfer.
\end{enumerate}

\textbf{What this branch book does not claim.}
\begin{itemize}
 \item The harmonization functor is not yet constructed.
 \item Matter content (quarks, leptons, Higgs) is not yet incorporated.
 \item The coupling constants $g_3, g_2, g_1$ are not yet derived.
 \item The three-generation mass hierarchy is a speculative analogy, not a theorem.
 \item No unification of gauge and gravitational forces is claimed.
\end{itemize}



%% ---------------------------------------------------------------
\section{Conditional SM Intrinsic Closure}\label{sec:sm-intrinsic-closure}
%% ---------------------------------------------------------------

\begin{proposition}[\status{C}]\label{prop:SM-intrinsic-closure}
\textup{(Conditional SM Intrinsic Closure.)}
\textup{[Conditional on the discharged/hardened forms of
O-SM.Harmonization (Ob.~\ref{ob:SM-harmonization}),
O-SM.Matter (Ob.~\ref{ob:SM-matter}),
O-SM.Higgs (Ob.~\ref{ob:SM-higgs}), and
O-SM.CouplingTransfer (Ob.~\ref{ob:SM-coupling-transfer-full}).]}
The SM super-branch has \emph{conditional intrinsic structural closure}:
its internal anti-smuggling obligations for gauge algebra, gauge-template
uniqueness, YM--GR interface compatibility, harmonization functor,
matter datum, Higgs/EWSB screening, and structural coupling transfer
are discharged at conditional force.

However, this proposition does \emph{not} promote the SM branch to
unconditional CERT-CLOSE. The SM super-branch:
\begin{enumerate}[label=\textup{(\arabic*)}]
 \item inherits all unresolved obligations of the YM and GR sub-branches;
 \item does not claim low-energy phenomenological matching, RG-running
 derivation, measured particle masses, CKM/PMNS data, or full
 empirical Standard Model reconstruction;
 \item does not claim unification of gauge and gravitational forces.
\end{enumerate}
\[
 \operatorname{Status}(\mathrm{SM}) \;\le\;
 \operatorname{Status}(\mathrm{YM}) \wedge
 \operatorname{Status}(\mathrm{GR}) \wedge
 \operatorname{Status}(\mathrm{SM}_{\mathrm{intrinsic}}).
\]
\end{proposition}

\begin{proof}
The four SM-internal obligations are conditionally discharged as stated:
O-SM.Harmonization by Thms.~\ref{thm:SM-harmonization-partial}
and~\ref{thm:F-SM-POT-morphism};
O-SM.Matter by Prop.~\ref{prop:rh-SM-matter-candidate}
and Thm.~\ref{thm:SM-matter-source-descended};
O-SM.Higgs by Prop.~\ref{prop:rh-SM-higgs-candidate}
and Thm.~\ref{thm:SM-higgs-source-descended};
O-SM.CouplingTransfer by Thm.~\ref{thm:SM-coupling-transfer-full}.
The SM super-branch is defined as a harmonization of YM and GR
endpoint data and explicitly inherits all open obligations of both
sub-branches by the super-branch governance rule; no SM-internal
discharge reduces those inherited obligations.
\qed
\end{proof}

\begin{proposition}[\status{C}]\label{prop:SM-status-reconciliation}
\textup{(SM Status Reconciliation.)}
Let SM-Old denote the status recorded in
Prop.~\ref{prop:SM-status} (CERT-PROJ with primary open obligations
O-SM.Harmonization, O-SM.Matter, O-SM.Higgs, O-SM.CouplingTransfer).
Let SM-New denote the later conditional hardening state in which those
four obligations are conditionally discharged at structural level.

Then SM-New supersedes SM-Old \emph{only with respect to} the
internal SM obligation ledger, not with respect to:
\begin{enumerate}[label=\textup{(\arabic*)}]
 \item inherited YM/GR obligation scope;
 \item CERT-CLOSE limitations (SM full CERT-CLOSE still unavailable);
 \item phenomenological non-claims (measured masses, CKM/PMNS,
 RG-running, full empirical reconstruction).
\end{enumerate}
The corrected status is:
\begin{center}
SM: \emph{conditionally intrinsic-structural closed,
inherited-scope open.}
\end{center}
This is strictly stronger than raw CERT-PROJ, but strictly weaker than
unconditional or full empirical CERT-CLOSE.

The following items in Prop.~\ref{prop:SM-status} are superseded
\emph{as internal ledger entries only}:
(i) ``The harmonization functor is not yet constructed'' ---
superseded in part by $F_{\mathrm{SM}}^{\mathrm{KPO}}$;
(ii) ``Electroweak breaking is still fully open'' ---
superseded at structural level by the $B_0$ screening discharge;
(iii) SM-three-gen was previously conditional; it has since been
upgraded to [U] by promotion of its former conditions.

The following warnings from Prop.~\ref{prop:SM-status} remain
fully binding: inherited YM/GR obligations, no unconditional
CERT-CLOSE, no phenomenological Standard Model completion.
\end{proposition}

\begin{proof}
Prop.~\ref{prop:SM-intrinsic-closure} gives the internal structural
closure; the inherited-scope and non-claims clauses follow directly
from the super-branch governance rule established in the preamble
the SM super-branch cannot be stronger than
the conjunction of its parent-branch assumptions, and the explicit
non-claims list is carried forward from the branch book boundary
the SM branch boundary declaration.
\qed
\end{proof}

\begin{remark}[\status{R}]\label{rem:sm-status-current}
\textbf{Current SM status after reconciliation.}
The most accurate current status line is:
\begin{center}
SM: \emph{conditionally intrinsic-structural closed, inherited-scope open.}
\end{center}
The next highest-leverage SM obligation is the extension of the
harmonization functor to the full matter sector:
$F_{\mathrm{SM}}^{\mathrm{KPO}+\mathcal{M}}$, requiring
O-SM.Matter and O-SM.Higgs at unconditional (rather than conditional
structural) force. That full canonical force requires the
O-ID\textsuperscript{cont} identification providing the full
continuum representation theory (see Thm.~\ref{thm:O-ID-cont-rate},
YM Branch).
\end{remark}

%% ---------------------------------------------------------------
\section{SM External Specification Frontier}
\label{sec:sm-extspec}
%% ---------------------------------------------------------------

The YM branch resolved its external specification gap via the
seven-predicate ClaySpec packet (YM Branch, §sec:ClaySpec). An
analogous specification question arises for the SM super-branch.
This section states the SM external specification predicate cleanly
and identifies the minimal remaining gaps.

\begin{definition}[\status{D}]\label{def:SM-extspec-predicate}
\textup{(SM External Specification Predicate.)}
The \emph{SM external specification predicate} holds when the
harmonized NFC SM object
$\mathcal{S}_{NFC} = (G_{\mathrm{SM}},\; \mathcal{H}_{\mathrm{phys}}^{SM},\;
H_{\mathrm{SM}},\; \mathfrak{h}_{\mathrm{SM}},\; \mathcal{M}_{\mathrm{SM}},\;
\phi_{B_0})$
satisfies:
\begin{enumerate}[label=\textup{(\arabic*)}]
 \item \emph{SM-G}: gauge algebra $\mathfrak{su}(3)\oplus
 \mathfrak{su}(2)\oplus\mathfrak{u}(1)$ (from O\_ID\textsuperscript{MSC});
 \item \emph{SM-HILB}: physical gauge-quotiented Hilbert space
 (from B1-NullID and SM harmonization);
 \item \emph{SM-MAT}: three-generation fermionic matter sector
 with correct quantum numbers and anomaly cancellation
 (from O-SM.Matter);
 \item \emph{SM-HIGGS}: $B_0$ Higgs screening projector realizing
 electroweak symmetry breaking to $\mathfrak{u}(1)_{\mathrm{em}}$
 (from O-SM.Higgs);
 \item \emph{SM-FORM}: harmonized Yang--Mills and gravitational
 action form on the common transport-stable core (from
 O-SM.CouplingTransfer);
 \item \emph{SM-OSW}: OS/Wightman structure at the SM harmonized
 scope (inherited from YM O\_CLU\textsuperscript{MSC}).
\end{enumerate}
\end{definition}

\begin{proposition}[\status{C}]\label{prop:SM-extspec-status}
\textup{(SM External Specification Status.)}
Under the current canon, the six SM-extspec predicates have the
following status:
\begin{enumerate}[label=\textup{(\arabic*)}]
 \item \emph{SM-G}: \textbf{[C]} --- O\_ID\textsuperscript{MSC}
 supplies $\mathfrak{su}(3)\oplus\mathfrak{su}(2)\oplus\mathfrak{u}(1)$
 at MSC-normalized scope.
 \item \emph{SM-HILB}: \textbf{[C]} --- B1-NullID and SM
 harmonization supply the physical gauge-quotiented Hilbert space.
 \item \emph{SM-MAT}: \textbf{[C] structural level} --- the
 three-generation quintuplet is source-descended at structural
 force (Prop.~\ref{prop:rh-SM-matter-candidate},
 Thm.~\ref{thm:SM-matter-source-descended});
 full \textbf{[U]} requires O\_ID\textsuperscript{cont} at
 unconditional force.
 \item \emph{SM-HIGGS}: \textbf{[C] structural level} --- the
 $B_0$ screening projection is source-descended at structural
 force (Prop.~\ref{prop:rh-SM-higgs-candidate},
 Thm.~\ref{thm:SM-higgs-source-descended});
 full \textbf{[U]} requires O\_ID\textsuperscript{cont} at
 unconditional force.
 \item \emph{SM-FORM}: \textbf{[C]} --- O-SM.CouplingTransfer
 discharged at structural level; the $6{:}2{:}1$ coupling ratio
 is source-descended.
 \item \emph{SM-OSW}: \textbf{[C]} --- inherited from YM
 O\_CLU\textsuperscript{MSC} at MSC-normalized scope.
\end{enumerate}
\textbf{SM external specification status:} \emph{SM-G, SM-HILB,
SM-FORM, SM-OSW conditionally closed; SM-MAT and SM-HIGGS
conditionally closed at structural level, blocked at unconditional
force by O\_ID\textsuperscript{cont}.}
\end{proposition}

\begin{proof}
Each item follows directly from the canonical references cited and
the intrinsic-structural closure proposition
(Prop.~\ref{prop:SM-intrinsic-closure}). The distinction between
structural force and unconditional force is governed by Book~VII
status discipline: SM-MAT and SM-HIGGS are proved from the declared
structural source data; unconditional force would require
O\_ID\textsuperscript{cont} providing the full continuum
representation theory.
\qed
\end{proof}

\begin{remark}[\status{R}]\label{rem:sm-extspec-frontier}
\textbf{SM minimal frontier: O\_ID\textsuperscript{cont}.}
The SM external specification predicate reduces to a single
unconditional gating dependency:
O\_ID\textsuperscript{cont} (the full identification of the
collar-local Lie algebra with $\mathfrak{su}(3)\oplus
\mathfrak{su}(2)\oplus\mathfrak{u}(1)$ at continuum representation
theory scope, providing unconditional force for the matter quantum
numbers and Higgs representation). All six SM-extspec predicates
are conditionally closed except the ``force upgrading'' needed for
SM-MAT and SM-HIGGS to reach unconditional level.
O\_ID\textsuperscript{cont} is not a new obligation introduced here;
it is already recorded as the rate-certificate theorem
(Thm.~\ref{thm:O-ID-cont-rate}, YM Branch) and its pending
extension to the full continuum representation theory.
\end{remark}

\end{document}
